\bookchapter{第十四章：パスワードマネージャー特集}

\chapterintro{一、なぜパスワードマネージャーが必要？}

\booksection{一、なぜパスワードマネージャーが必要？}

パスワードマネージャーがこのジレンマを解決：マスターパスワード一つだけ覚えれば、他のすべてのパスワードを覚えてくれます。

二、パスワードマネージャーの仕組み

\booksection{二、パスワードマネージャーの仕組み}

\begin{enumerate}
  \item \textbf{：すべてのパスワードは高強度暗号化で保存され、あなただけが復号可能。}マスターパスワード
  \item \textbf{：パスワード保管庫をロック解除するためのマスターパスワードを設定。マスターパスワード自体はサーバーに保存されません。}自動入力
  \item \textbf{：ウェブサイトやAppのログイン画面で、パスワードマネージャーが自動的にIDとパスワードを入力。}パスワード生成
  \item \textbf{：ランダムで複雑なパスワードを生成し、セキュリティを向上。}三、主流パスワードマネージャーのおすすめ
\end{enumerate}

\booksection{三、暗号化原理の概要}

パスワードマネージャーの暗号化原理を理解すると、サービスプロバイダーが信頼に値するかを判断できます。以下、分かりやすい言葉でいくつかの核心概念を説明します。

\booksubsection{AES-256暗号化}

AES（Advanced Encryption Standard）は米国国立標準技術研究所（NIST）が認可した共通鍵暗号アルゴリズムで、銀行、軍、政府で広く使用されています。「256」は鍵長が256ビットであることを意味し——理論上、宇宙の年齢を超える時間がないと総当たり攻撃で破れません。主流のパスワードマネージャー（1Password、Bitwarden、KeePassなど）はすべてAES-256でパスワード保管庫を暗号化しています。

\booksubsection{ゼロナレッジアーキテクチャ（Zero-Knowledge Architecture）}

「ゼロナレッジ」とはサービスプロバイダーが何も保存しないということではなく、\textbf{暗号化と復号はあなたのデバイスのローカルで完結し、サーバーは暗号化済みの暗号文のみを保存する}。具体的な流れ：

\begin{enumerate}
  \item クライアントでマスターパスワードを入力。
  \item クライアントはマスターパスワードから暗号化鍵を派生させます（この処理はあなたのスマホ/PC上で行われます）。
  \item 暗号化された暗号文のみがサーバーにアップロードされます。
  \item 次回ログイン時、サーバーは暗号文を送信し、クライアントが再びマスターパスワードで復号します。
\end{enumerate}

重要なポイント：\textbf{マスターパスワード自体はあなたのデバイスから一切送信されません}。サーバー上にもありません。そのため、仮にサービス提供元がハッキングされても、攻撃者が手にするのは解読不可能な暗号文の山に過ぎません——マスターパスワードを推測されない限り。

\begin{tcbinfo}{なぜ「ゼロナレッジ」はサービス提供者にもパスワードを見せないのか？}
暗号化鍵はあなたのマスターパスワードから派生し、マスターパスワードはあなたの頭の中にのみ存在し、ローカルデバイス上でのみ使用されるからです。サービス提供者のサーバーにはマスターパスワードも暗号化鍵もなく、暗号文しかありません。だからこそ、Bitwarden、1Password などの正規サービス提供者は「私たちもあなたのパスワードを見ることはできません」と明確に宣言しているのです——見たくないのではなく、技術的に不可能なのです。

\end{tcbinfo}

\booksubsection{鍵派生関数：PBKDF2 と Argon2}

マスターパスワードは暗号化鍵としてそのまま使えません——「鍵派生関数」による処理が必要です。この関数の役割は\textbf{あえて計算を遅くすること}ことで、ハッカーの総当たり攻撃のコストを非現実的なレベルまで引き上げます。

\begin{itemize}
  \item \textbf{PBKDF2}: 古い標準規格で、iCloud キーチェーン、Bitwarden、LastPass などに広く使われています。反復ハッシュ（通常10万回以上）で破解を遅延させます。
  \item \textbf{Argon2}: 2015年パスワードハッシュコンテスト優勝で、PBKDF2 よりモダンです。「計算コスト」に加え「メモリコスト」も導入し、GPU/ASIC などの専用ハードウェアでも加速破解が困難です。Bitwarden、KeePassXC などが既にサポートしています。
\end{itemize}

反復回数が多いほど安全ですが、ロック解除速度も遅くなります。正規サービス提供者は使用する派生アルゴリズムと反復回数を公開し、セキュリティコミュニティによる監査を可能にしています。

\booksubsection{事例警鐘：LastPass 2022年データ漏洩}

2022年8月、老舗パスワードマネージャー LastPass は開発環境への侵入を公表しました。同年11月にはさらに、ハッカーがユーザーのパスワード保管庫のバックアップ（暗号化された暗号文）を持ち出したことを明らかにしました。この事件の重要な教訓：

\begin{itemize}
  \item \textbf{ゼロナレッジアーキテクチャが最低限を守った}: パスワード保管庫自体が暗号化されているため、ハッカーは暗号文を手にしても解読できません——ユーザーのマスターパスワードが十分に強いことが前提です。
  \item \textbf{弱いマスターパスワードのユーザーが最も危険}: ハッカーは弱いマスターパスワードをオフラインで総当たり攻撃できるため、弱いマスターパスワードや使い回しているユーザーは実質的なリスクに直面します。
  \item \textbf{メタデータも漏洩}: パスワードの内容は暗号化されていますが、サイトURLやメモなどのメタデータは暗号化されていない可能性があり、ユーザーがどのサービスを利用しているかが露出します。
\end{itemize}

\begin{tcbwarn}{LastPass 事件から学んだこと}
第一に、マスターパスワードは長くユニークである必要があります（16文字以上、大文字·小文字·数字·記号を含む）。第二に、Argon2id をサポートし反復回数の多いサービス提供者を優先しましょう。第三に、高価値アカウント（銀行、メール）には必ず独立した二要素認証を設定し、パスワードマネージャーだけに頼らないようにしましょう。

\end{tcbwarn}

\booksection{四、主要パスワードマネージャーのおすすめ}

\booksubsection{iCloudキーチェーン}

\begin{itemize}
  \item \textbf{：Appleエコシステムユーザーの第一選択。Safari、Appとシームレスに統合、無料、安全。欠点はクロスプラットフォーム利用が不便。}Googleパスワードマネージャー
  \item \textbf{：Androidユーザーに便利。Chrome、Androidシステムと統合、無料。}Huaweiパスワード金庫
  \item \textbf{：HarmonyOS内蔵、Huaweiユーザーに適しています。}サードパーティ専門ツール：
\end{itemize}

\booksubsection{Bitwarden}

\begin{itemize}
  \item \textbf{：オープンソース無料、機能強力、すべてのプラットフォームに対応。プライバシー重視ユーザーに強くおすすめ。}1Password
  \item \textbf{：インターフェース美麗、機能豊富、家庭やチームに適しています。有料。}LastPass
  \item \textbf{：老舗パスワードマネージャー、無料版は機能制限あり。}KeePass
  \item \textbf{：完全ローカル保存、オープンソース無料。ギークユーザーに適しています。}四、パスワードマネージャー活用のベストプラクティス
\end{itemize}

\booksection{五、パスワード強度の本当の評価}

パスワードマネージャーは一般的に「パスワード強度」インジケーター（赤/黄/緑）を表示しますが、この評価は長さと文字の種類しか見ておらず、包括的ではありません。パスワードが本当に安全かを判断するには「エントロピー」と「漏洩の有無」を考慮する必要があります。

\booksubsection{エントロピー（Entropy）：パスワード強度の科学的尺度}

「エントロピー」はパスワードの不確実性を測る指標で、単位はビットです。エントロピーが高いほど、総当たり攻撃に必要な試行回数が多くなります。

\begin{itemize}
  \item \textbf{"123456"}: エントロピーが極めて低く、瞬時に破解されます。
  \item \textbf{"P@ssw0rd"}: 一見複雑に見えますが、すべての辞書攻撃リストに載っており、エントロピーは依然として低いです。
  \item \textbf{"correct horse battery staple"}: 4つのランダムな単語、エントロピー約44ビットで、破解難度は「P@ssw0rd」をはるかに超えます。
  \item \textbf{16文字のランダム文字（大文字·小文字+数字+記号を含む）}: エントロピー約100ビット、現在の計算力では総当たり攻撃不可能です。
\end{itemize}

結論：\textbf{「長さ + ランダム性」は「複雑な記号」よりも重要}。これがパスワードマネージャーが生成するランダムな長いパスワードが本当に安全である理由です。

\booksubsection{HaveIBeenPwned：あなたのパスワードが漏洩していないか確認}

セキュリティ研究者 Troy Hunt が運営する\href{https://haveibeenpwned.com}{haveibeenpwned.com}は過去の大規模データ漏洩で流出したアカウントとパスワードを収録しています。以下のことができます：

\begin{enumerate}
  \item メールアドレスを入力して、あなたのアカウントがいずれかの漏洩事件に関与していないか確認。
  \item Pwned Passwords 機能を使って、特定のパスワードが既知の漏洩データベースに含まれているか確認（パスワードマネージャーやブラウザにこの機能が組み込まれています）。
\end{enumerate}

もしパスワードが漏洩データベースに含まれていたら、\textbf{どんなに複雑に見えてもすぐに変更する必要があります}——辞書攻撃はこれらの既知パスワードを優先的に試すからです。

\begin{tcbinfo}{k-匿名性クエリの原理}
HaveIBeenPwned にパスワードを入力するとまた漏洩するのでは？心配いりません。現代のクエリは「k-匿名」プロトコルを採用しています：ブラウザはパスワードハッシュの最初の5文字だけをサーバーに送信し、サーバーはその5文字で始まる全ハッシュリストを返し、ローカルで照合します。つまり、サーバーはあなたの完全なパスワードハッシュを見ることはありません。1Password、Bitwarden の「漏洩パスワード検出」もこの仕組みで動作しています。

\end{tcbinfo}

\booksubsection{パスワードマネージャー内蔵の監視機能}

\begin{itemize}
  \item \textbf{1Password Watchtower}: 弱いパスワード、重複パスワード、漏洩パスワード、期限切れパスワード、2FA未設定サイトをリアルタイムで監視し、赤/オレンジ/黄のレベル分けで表示します。
  \item \textbf{Bitwarden Vault Health}: 無料版でも弱いパスワード、重複パスワードを確認可能。有料版では漏洩パスワードの検出が可能です（HaveIBeenPwned ベース）。
  \item \textbf{iCloud キーチェーン「セキュリティレコメンデーション」}: iOS 設定 パスワード セキュリティレコメンデーションで、どのパスワードが既知の漏洩に含まれているか、どれが重複使用されているかを教えてくれます。
  \item \textbf{Google パスワードマネージャー「パスワードチェックアップ」}: Chrome に内蔵され、漏洩したパスワード、弱いパスワード、重複パスワードをワンクリックでチェックできます。
\end{itemize}

\begin{tcbpractice}{今すぐ確認：あなたのパスワードは安全か}
\begin{enumerate}
  \item haveibeenpwned.com を開き、よく使うメールアドレスを入力して、どの漏洩事件に関与しているか確認しましょう。
  \item iPhone ユーザー：設定 パスワード セキュリティレコメンデーションで、赤色の警告を一つずつ処理しましょう。
  \item Chrome ユーザー：設定 パスワード パスワードチェックアップで、指示に従って変更しましょう。
  \item 1Password/Bitwarden ユーザー：Watchtower/Health Reports を開き、「漏洩」と「重複」の項目を優先的に処理しましょう。
\end{enumerate}

\end{tcbpractice}

\booksection{六、Passkey：パスワード不要ログインの未来}

2024-2026年、パスワード技術は革命を遂げようとしています——\textbf{Passkey（パスキー）}が従来のパスワードを完全に淘汰する可能性があります。Apple、Google、Microsoft の3大プラットフォームが共同推進し、FIDO Alliance と W3C が支持する、数十年で最大のログイン方式の変革です。

\booksubsection{Passkey とは？}

Passkey は FIDO2 / WebAuthn 標準に基づき、\textbf{デバイス上の鍵ペア}でパスワードを置き換えます。原理を簡略化すると：

\begin{itemize}
  \item あなたのスマホが鍵のペアを生成します：\textbf{公開鍵}（ウェブサイトに登録）+\textbf{秘密鍵}（デバイスのセキュリティチップに永久保存）。
  \item ログイン時、ウェブサイトが「チャレンジ」を送信し、あなたのスマホが秘密鍵で署名し、ウェブサイトが公開鍵で検証します。
  \item 署名前に生体認証（Face ID / 指紋）またはデバイスの PIN で本人確認が必要です。
  \item このプロセス全体で、\textbf{パスワードは一切ネットワーク上で送信されません}。
\end{itemize}

\booksubsection{Passkey vs 従来のパスワード}

\begin{center}
\begin{tabularx}{\linewidth}{|X|X|X|}
\hline
\textbf{比較項目} & \textbf{従来のパスワード} & \textbf{Passkey} \tabularnewline
\hline
フィッシング耐性 & 低い、フィッシングサイトがパスワードを盗める & 本質的にフィッシング耐性、秘密鍵は本物のドメインにのみ署名 \tabularnewline
サーバー漏洩の影響 & パスワードハッシュが破解される可能性がある & サーバーは公開鍵のみ保存、漏洩の影響なし \tabularnewline
利便性 & 記憶/入力/管理が必要 & Face ID/指紋で1秒ログイン \tabularnewline
リプレイ攻撃 & 中間者攻撃で傍受·再生される可能性 & 毎回使い捨て署名、再生不可 \tabularnewline
\hline
\end{tabularx}
\end{center}

\booksubsection{各プラットフォームのサポート状況}

\begin{itemize}
  \item \textbf{Apple iCloud キーチェーン Passkey}: iOS 16+ / macOS Ventura+ でネイティブサポート。Passkey は同じ Apple ID でログインしている全デバイス間で同期でき、新デバイスでサイトにログインする際に再作成不要です。
  \item \textbf{Google パスワードマネージャー Passkey}: Android + Chrome でネイティブサポート。Passkey は Google アカウントに保存され、Android デバイス間で同期。2023年以降、Google 個人アカウントで Passkey ログインをサポート。
  \item \textbf{Microsoft}: Windows 11 / Edge が Passkey をサポート。Microsoft アカウント、Outlook、Office などで順次 Passkey ログインを提供し、Microsoft Authenticator で管理可能。
  \item \textbf{1Password / Bitwarden}: Passkey の作成、保存、クロスプラットフォーム同期をサポートし、単一エコシステムに縛られません。
\end{itemize}

\begin{tcbinfo}{クロスデバイス同期が重要なブレークスルー}
初期の Passkey の最大の課題は「単一デバイスにバインド」——スマホを替えるたびに再作成が必要でした。2023年以降、Apple、Google、Microsoft は Passkey のクラウド同期を実現：iCloud キーチェーンは Apple デバイス間で同期、Google パスワードマネージャーは Android デバイス間で同期。1Password などのサードパーティはクロスエコシステム同期を実現しました。これは Passkey がパスワードを代替する実用性を備えたことを意味します。

\end{tcbinfo}

\booksubsection{スマホで Passkey を作成·使用する方法}

\textbf{iOS（iPhone）の操作手順：}

\begin{enumerate}
  \item Passkey をサポートするウェブサイト（Google、Microsoft、Amazon、eBay、PayPal など）のアカウントセキュリティページで「パスキーを追加 / Add passkey」を選択。
  \item システムが Face ID / Touch ID の認証画面を表示し、認証成功後 Passkey が自動生成され iCloud キーチェーンに保存されます。
  \item 次回そのサイトにログインする際、ユーザー名入力後に「iPhone でログイン」を選択し、Face ID を一度認証するだけです。
  \item 管理画面：設定 パスワード 該当サイトを見つけると Passkey エントリーが表示されます。
\end{enumerate}

\textbf{Android の操作手順：}

\begin{enumerate}
  \item Passkey をサポートするウェブサイトのアカウントセキュリティページで「パスキーを追加」を選択。
  \item Google パスワードマネージャーのポップアップ Google アカウントを選択 画面ロック認証（指紋/パターン/PIN）。
  \item Passkey は Google パスワードマネージャーに自動保存され、同じアカウントの Android デバイス間で同期されます。
  \item 次回ログイン時に「Passkey を使用」を選択し、画面ロックを解除するだけです。
  \item 管理画面：Android 設定 パスワード、パスキー、自動入力 パスワードマネージャー google.com。
\end{enumerate}

\booksubsection{Passkey のクロスデバイスログイン}

Windows PC で、iPhone にしか Passkey がないサイトにログインする場合、QRコードで認証できます：

\begin{enumerate}
  \item サイトのログインページで「Passkey を使用」を選択。
  \item 「他のデバイスを使用」を選択。
  \item PC にQRコードが表示されるので、iPhone でスキャン。
  \item iPhone に Face ID 認証が表示され、確認後 PC 上でサイトのログインが完了します。
\end{enumerate}

このプロセスは Bluetooth + インターネットの二重暗号化通信で行われ、Passkey の秘密鍵が iPhone を離れることはありません。

\begin{tcbpractice}{今すぐ試す：パスワード不要ログインを体験}
\begin{enumerate}
  \item Google アカウントにログイン セキュリティ パスキー パスキーを追加で、Face ID/指紋ログインを体験。
  \item Microsoft アカウントにログイン セキュリティ 詳細セキュリティオプション パスキーを追加。
  \item タオバオ、JD（京東）、eBay、PayPal など Passkey 対応サービスで試してみましょう。
  \item 従来のパスワードをバックアップとして残し、高価値アカウントを徐々に Passkey に移行しましょう。
\end{enumerate}

\end{tcbpractice}

\begin{tcbwarn}{移行期のアドバイス}
Passkey はパスワードを完全に置き換えたわけではありません——多くのサイトはまだ「Passkey + パスワード」の併用段階です。高価値アカウントで Passkey を優先的に有効化しますが、\textbf{従来のパスワードをすぐに削除しないでください}、バックアップとして残しましょう。同時にパスワードマネージャーの使用も続けてください。すべてのサイトが Passkey に対応しているわけではないからです。Passkey は「パスワードマネージャーの強力な補完」として捉え、「即時置き換え」とは考えないでください。

\end{tcbwarn}

\booksection{七、ハードウェアセキュリティキー}

一般ユーザーより高いセキュリティを求める場合——例えば大量の暗号資産を保有している、経営者/ジャーナリスト/活動家である、企業のコアアカウントを保護するなど——ハードウェアセキュリティキーはSMS認証コードより信頼性の高い選択肢です。

\booksubsection{ハードウェアセキュリティキーとは？}

ハードウェアキーは物理デバイス（USBメモリのような形状）で、内部のセキュリティチップにあなたの秘密鍵を保存し、物理的接触または NFC で認証を完了します。代表的な製品：

\begin{itemize}
  \item \textbf{YubiKey}（Yubico社）：最も主流。FIDO2/WebAuthn、U2F、OTP、PIV、OpenPGP をサポートし、USB-C / USB-A / NFC など多様なインターフェースに対応。5シリーズで50-70ドル。
  \item \textbf{Google Titan}: Google 製。USB-C + NFC / Bluetooth、約30-40ドル。
  \item \textbf{OnlyKey}: オープンソースの代替品、複数プロトコルをサポート。
  \item \textbf{中国産代替品}: 飛天誠信 ePasskey など、より手頃な価格。
\end{itemize}

\booksubsection{なぜハードウェアキーはSMS認証コードより安全なのか？}

\begin{itemize}
  \item \textbf{フィッシング耐性}: Passkey と同様、ハードウェアキーは本物のドメインにのみ応答し、フィッシングサイトは認証コードを騙し取れません。
  \item \textbf{SIM ハイジャック耐性}: SMS認証コードはSIMスワップ攻撃で傍受される可能性がありますが、ハードウェアキーは電話番号を必要としません。
  \item \textbf{中間者攻撃耐性}: FIDO2 プロトコルは毎回の署名にチャレンジ値を含み、リプレイできません。
  \item \textbf{物理的分離}: 秘密鍵は常にハードウェアチップ内にあり、マルウェアがPCを侵害しても読み取れません。
\end{itemize}

\booksubsection{適用シーン}

\begin{itemize}
  \item 暗号資産取引所アカウント、ウォレット復元。
  \item 企業SSO、リモートワークVPN、SSHサーバーログイン。
  \item メール、GitHub、ソーシャルメディアなどの「メインアカウント」——一度侵害されると他の全アカウントがリセットされます。
  \item ジャーナリスト、弁護士、経営者など高感度人物。
\end{itemize}

\booksubsection{ハードウェアキーを二要素認証に設定する方法}

Google アカウント + YubiKey を例に：

\begin{enumerate}
  \item YubiKey を購入（2本まとめて購入を推奨、1本は日常用、1本はバックアップ、別々の場所に保管）。
  \item myaccount.google.com にログイン セキュリティ 2段階認証。
  \item 「セキュリティキーを追加」を選択 YubiKey を USB-C ポートに挿入、またはスマホのNFCセンサーにかざす。
  \item 指示に従って YubiKey の金属接点に触れて登録を完了。
  \item 2本目のバックアップキーの登録を推奨。
  \item 「2段階認証」の設定で、ハードウェアキーを優先認証方式にし、SMS認証コードをバックアップにします。
\end{enumerate}

主流サポートサイト：Google、GitHub、Microsoft、Apple ID、Facebook、Twitter、Dropbox、GitLab、AWS、Coinbase、Binance など。

\begin{tcbinfo}{デュアルキー戦略}
ハードウェアキーの最大のリスクは「紛失」です。推奨：(1) 異なる2本のキーを同時に登録；(2) 1本は随身携帯（キーホルダー）、1本は自宅の金庫または銀行貸金庫に保管；(3) 一部サイトでは「バックアップ復旧コード」を印刷してオフライン保存可能。これで主キーを紛失してもアカウントから永久に締め出されることはありません。

\end{tcbinfo}

\begin{tcbwarn}{購入ルートに注意}
必ず公式ストアまたは正規代理店から購入し、中古や出所不明のキーは買わないでください——理論上、プリセットされたハードウェアにバックドアが仕込まれている可能性があります。YubiKey のファームウェアは出荷後に更新できませんが、Yubico は継続的にセキュリティ公告を発信するので公式サイトを確認してください。

\end{tcbwarn}

\booksection{八、家族/チームのパスワード共有}

家庭では Netflix、家庭用 Wi-Fi、ルーター管理画面、スマートロックのパスワードを家族で共有。チームではサーバー、データベース、ソーシャルメディアアカウントを複数人で共同利用します。\textbf{WeChat でパスワードを送信したり、共有ドキュメントに書いたりするのは高危険}——チャット履歴やドキュメント権限が漏洩すれば、パスワードも漏洩します。

\booksubsection{1Password の「共有保管庫」（Vault）}

1Password は家族·チーム共有の代表的な製品です。その核心仕組みは\textbf{複数の保管庫}：

\begin{itemize}
  \item \textbf{個人保管庫}: あなただけが閲覧可能。銀行、メールなどのプライベートアカウント。
  \item \textbf{共有保管庫}: 家族やチームメンバー全員が閲覧可能。Netflix、家庭用 Wi-Fi、共有クレジットカードなど。
  \item \textbf{権限レベル分け}: チーム版では一部メンバーを「閲覧のみ」や「閲覧のみ·編集不可」に設定できます。
\end{itemize}

家族やチームメンバーはマスターパスワードを共有する必要がなく、各自のアカウントでログインしますが、権限を付与された共有保管庫の内容を閲覧できます。パスワード変更後、全メンバーにリアルタイムで同期されます。

\booksubsection{Bitwarden の「組織」機能}

Bitwarden の無料版は個人パスワード共有をサポート（「Send」機能で期限付きで1件のパスワードを共有）。有料版では「組織」を作成し、複数エントリーを家族やチームメンバーと共有でき、権限を細かく設定可能。Bitwarden はセルフホストオプションも提供し、データ主権を重視する企業に適しています。

\booksubsection{家庭シーン：Netflix、Wi-Fi などのパスワードを安全に共有する方法}

\begin{itemize}
  \item \textbf{Netflix / 動画サイト}: 1Password Families の「共有保管庫」にログインパスワードを保存し、家族メンバーが各自でログイン。グループチャットに送信する必要なし。
  \item \textbf{家庭用 Wi-Fi パスワード}: iOS デバイス間では AirDrop で来客に Wi-Fi パスワードを直接共有でき、口頭で伝える必要なし。Android 12+ にも「近距離共有でWi-Fi共有」機能があります。
  \item \textbf{スマートロック、スマート家電}: なるべくメーカーAppの「家族メンバー」機能で各自に個別アカウントを作成し、1つの管理者アカウントを共有しないようにしましょう。
  \item \textbf{家庭用ルーター管理画面}: 管理者パスワードは 1Password 共有保管庫に保存し、ゲスト用 Wi-Fi は別の SSID で設定。
\end{itemize}

\begin{tcbinfo}{1Password Families のコストパフォーマンス}
1Password Families は約4.99ドル/月で最大5人の家族メンバーをサポート。各人に独立した個人保管庫 + 4つの共有保管庫。一人当たり1ドル/月未満で、現在家庭共有で最も優れた体験を提供します。Bitwarden の家庭版は10ドル/年で6人サポートし、コストパフォーマンスはより高いですが、インターフェース体験はやや劣ります。

\end{tcbinfo}

\booksubsection{チームシーン：サーバー/データベースパスワード共有}

\begin{itemize}
  \item \textbf{1Password Teams / Business}: 部門やプロジェクトグループごとに複数の保管庫を作成し、きめ細かい権限管理が可能。すべてのアクセス行為を監査できます。
  \item \textbf{Bitwarden Teams / Enterprise}: オープンソース方案、セルフホスト可能。エンタープライズ版はSSO連携、ディレクトリ同期をサポート。
  \item \textbf{職務分離の原則}: 管理者パスワードを一人だけが掌握すべきではありません——2-3名の「緊急連絡先」を設定し、キーパーソンの離職/休暇時も業務が継続するように。
  \item \textbf{離職プロセス}: 従業員の離職時にパスワードマネージャーのアクセス権を即時取り消し、必要に応じてその従業員が接触した高感度パスワードをすべてローテーション。
\end{itemize}

\booksubsection{緊急アクセス（Emergency Access）：万が一の時のために}

こんなシーンを想像してください：あなたが家庭のすべての重要アカウントを単独で管理しているが、意図せず入院したり亡くなったりした場合、家族が銀行、メール、保険のアカウントにログインできない。\textbf{「緊急アクセス」機能}はまさにこのために設計されました——信頼する家族/友人を認可し、あなたが応答できない時にパスワード保管庫にアクセスできるようにします。

\textbf{Bitwarden 緊急アクセスの設定：}

\begin{enumerate}
  \item Bitwarden Web Vault にログイン 設定 緊急アクセス。
  \item 信頼する連絡先のメールアドレスを追加し、相手を招待。
  \item 「待機期間」を設定（7-14日を推奨）：相手がリクエストしてからあなたが確認するまでの時間枠。
  \item 待機期間内にあなたが応答しない（拒否もしない）場合、期間満了後に相手が自動的にアクセス権を取得します。
  \item いつでも手動でリクエストを拒否したり緊急アクセス権を取り消すことができます。
\end{enumerate}

\textbf{1Password 緊急アクセス：}

\begin{enumerate}
  \item 1Password は直接「緊急アクセス」とは呼びませんが、「家族/チーム組織」機能で実現できます——家族を家族版に追加した後、一部エントリーを「共有保管庫」に配置。
  \item 個人版ユーザーには、マスターパスワードを高度に信頼する家族に伝えるか、密封した封筒で金庫に保管することを推奨。
  \item 1Password は「復旧コード」も提供しており、二要素認証を有効にする際は必ず印刷して保管してください。
\end{enumerate}

\begin{tcbpractice}{今すぐ設定：家族のために「スペアキー」を残す}
\begin{enumerate}
  \item 信頼する家族1人を「緊急連絡先」に選定。
  \item パスワードマネージャーで緊急アクセスを設定（Bitwarden）するか、家族共有保管庫に追加（1Password）。
  \item マスターパスワード + 復旧コードを紙に書き、封筒に入れて自宅の金庫に保管するか信頼できる家族に伝える。
  \item 半年に一度見直し、緊急連絡先が引き続き適切か確認。
\end{enumerate}

\end{tcbpractice}

\booksection{九、パスワードマネージャー活用のベストプラクティス}

\begin{itemize}
  \item \textbf{：マスターパスワードは最後の砦、必ず長く複雑なパスワードを設定し、しっかり記憶。}二要素認証を有効化
  \item \textbf{：パスワードマネージャーアカウントで二要素認証をオンにし、保護を強化。}定期的に弱パスワードをチェック
  \item \textbf{：ほとんどのパスワードマネージャーには「セキュリティ監査」機能があり、弱パスワードや重複パスワードを見つけ出します。}マスターパスワードを忘れない
  \item \textbf{：マスターパスワードは復元不可、忘れるとすべてのパスワードを失います。紙に書いて安全な場所に保管推奨。}緊急アクセス
  \item \textbf{：緊急連絡先を設定、万が一アクセスできなくなった場合、信頼できる人があなたのパスワードを取得可能。}重要な注意
\end{itemize}

\begin{tcbwarn}{パスワードマネージャーは便利ですがリスクもあります：マスターパスワードが漏洩すると、すべてのパスワードが露出します。だからマスターパスワードは他のパスワードと絶対に同じにせず、見つけやすい場所に書かないでください。また、良好なセキュリティ実績のあるサービスプロバイダーを選ぶことをおすすめします。}
今すぐ試す：パスワードマネージャーを使い始める

\end{tcbwarn}

\begin{tcbpractice}{パスワードマネージャーを一つ選ぶ（Bitwardenまたはシステム標準を推奨）。}
\begin{enumerate}
  \item アカウント登録、強力なマスターパスワードを設定。
  \item 既存のパスワードをインポート（あれば）または新規パスワードの追加を開始。
  \item ブラウザプラグインまたはスマホAppをインストールし、自動入力機能を体験。
  \item セキュリティ監査レポートを確認し、改善が必要なパスワードを見つける。
  \item 三、応用編（第15-20章）
\end{enumerate}

\end{tcbpractice}
